\chapter{Présentation de Reminex}
\label{chap:reminex}

Reminex est l'entité au sein de laquelle j'ai effectué mon stage. Les informations de ce chapitre s'appuient principalement sur la documentation de présentation de l'entreprise (\textit{Expertise \& Capabilities})~\cite{reminex_capabilities} et sur les explications qui m'ont été données pendant le stage.

\section{La filiale d'ingénierie du Groupe Managem}

Reminex est la filiale d'ingénierie et de recherche du Groupe Managem. Sa vocation est d'accompagner les projets miniers du Groupe --- et de clients extérieurs --- depuis l'exploration jusqu'à la mise en production, puis au-delà, pendant l'exploitation. Elle réunit, d'après sa documentation, plus de quatre-vingts spécialistes et revendique plusieurs dizaines d'années d'expérience cumulée en projets et en opérations minières, une dizaine de brevets et un large portefeuille de procédés développés en interne~\cite{reminex_capabilities}.

Ce positionnement est particulier~: parce qu'elle appartient à un groupe qui exploite lui-même des mines, Reminex combine une culture d'ingénierie « de bureau » et une culture d'exploitation « de terrain ». C'est un point qui revient souvent dans sa communication et que les ingénieurs m'ont confirmé~: concevoir en connaissant les contraintes réelles de l'exploitation.

\section{Trois missions : créer, livrer, améliorer}

L'entreprise résume sa mission autour de trois verbes, qui structurent aussi la manière dont elle intervient sur un projet~\cite{reminex_capabilities}~:

\begin{center}
\begin{tikzpicture}[
  node distance=6mm,
  mission/.style={rectangle, rounded corners, draw=reminexorange!80, fill=reminexorange!12,
    text width=3.3cm, minimum height=1.6cm, align=center, font=\small},
  arr/.style={-{Stealth[length=3mm]}, reminexorange!80, line width=1pt}]
\node[mission] (c) {\textbf{Créer}\\[2pt]\scriptsize concevoir des solutions et des \textit{flowsheets} pour valoriser les gisements};
\node[mission, right=of c] (d) {\textbf{Livrer}\\[2pt]\scriptsize conception, ingénierie et exécution du projet jusqu'à sa réussite};
\node[mission, right=of d] (i) {\textbf{Améliorer}\\[2pt]\scriptsize accompagner l'exploitation pour améliorer durablement l'actif};
\draw[arr] (c) -- (d);
\draw[arr] (d) -- (i);
\end{tikzpicture}
\captionof{figure}{Les trois missions de Reminex (d'après la documentation de l'entreprise~\cite{reminex_capabilities}).}
\label{fig:missions}
\end{center}

Autrement dit, l'entreprise ne s'arrête pas à la remise des plans~: elle intervient en amont (études, conception du procédé) et en aval (mise en service, optimisation de l'usine en production).

\section{Un éventail de disciplines d'ingénierie}

La conception d'une usine minière mobilise simultanément de nombreuses spécialités. Reminex dispose en interne d'une équipe pluridisciplinaire couvrant notamment les disciplines suivantes~\cite{reminex_capabilities}~:

\begin{center}
\renewcommand{\arraystretch}{1.25}
\begin{tabular}{@{}p{3.4cm} p{10.5cm}@{}}
\toprule
\textbf{Discipline} & \textbf{Rôle dans le projet} \\
\midrule
Procédé (\textit{Process}) & Définit le \textit{flowsheet}, les bilans matière et les conditions de fonctionnement. \\
Mécanique & Choix et dimensionnement des équipements (concasseurs, pompes, convoyeurs\ldots). \\
Génie civil & Fondations, ouvrages en béton, terrassements. \\
Structure & Charpentes métalliques supportant les équipements. \\
Tuyauterie (\textit{Piping}) & Réseaux de fluides reliant les équipements. \\
Instrumentation & Capteurs et boucles de mesure/régulation. \\
Électricité & Distribution et motorisation. \\
Environnement & Gestion des rejets, de l'eau, des stériles. \\
Gestion de projet & Planning, coûts, coordination des disciplines. \\
\bottomrule
\end{tabular}
\captionof{table}{Principales disciplines d'ingénierie de Reminex.}
\label{tab:disciplines}
\end{center}

La particularité, que j'ai découverte pendant le stage, est que ces disciplines ne travaillent pas les unes après les autres mais \emph{ensemble}~: un choix de procédé impose une structure, qui impose une tuyauterie, qui impose une alimentation électrique, etc. La coordination entre elles est un métier à part entière (la gestion de projet).

\section{Un accompagnement « de la mine au produit »}

Enfin, Reminex décrit son offre comme un accompagnement continu couvrant tout le cycle d'un projet minier, « de la fosse au produit »~\cite{reminex_capabilities}~: exploration et caractérisation, développement du \textit{flowsheet}, ingénierie et conception, gestion de la construction (jusqu'à l'\sig{epcm}), mise en service et démarrage, puis excellence opérationnelle pendant la production. Ce continuum est important pour comprendre la suite du rapport~: chaque projet du Groupe se situe à un moment précis de ce cycle.
